\documentclass[10pt, letterpaper]{article}

\usepackage[margin=0.72in]{geometry}
\usepackage{array}
\usepackage{booktabs}
\usepackage{longtable}
\usepackage{tabularx}
\usepackage{pdflscape}
\usepackage[table]{xcolor}
\usepackage{titlesec}
\usepackage{fancyhdr}
\usepackage{enumitem}
\usepackage{microtype}
\usepackage[hidelinks]{hyperref}

\definecolor{headingblue}{HTML}{1F4E79}
\definecolor{lightblue}{HTML}{D9EAF7}
\definecolor{softgray}{HTML}{F2F4F7}
\definecolor{linegray}{HTML}{A6A6A6}

\pagestyle{fancy}
\fancyhf{}
\lhead{\small Md Kamruzzaman Kamrul}
\rhead{\small EDCI 59100-016 \textbar{} Summer 2026}
\cfoot{\small \thepage}
\renewcommand{\headrulewidth}{0.4pt}
\setlength{\headheight}{14pt}

\titleformat{\section}{\normalfont\large\bfseries\color{headingblue}}{\thesection.}{0.45em}{}
\titleformat{\subsection}{\normalfont\normalsize\bfseries\color{headingblue}}{\thesubsection.}{0.45em}{}
\titlespacing*{\section}{0pt}{10pt}{4pt}
\titlespacing*{\subsection}{0pt}{7pt}{3pt}

\setlength{\parindent}{0pt}
\setlength{\parskip}{4pt}
\setlist[itemize]{leftmargin=*, itemsep=2pt, topsep=2pt}
\setlist[enumerate]{leftmargin=*, itemsep=2pt, topsep=2pt}
\renewcommand{\arraystretch}{1.16}

\newcolumntype{L}[1]{>{\raggedright\arraybackslash}p{#1}}
\newcolumntype{Y}{>{\raggedright\arraybackslash}X}

\begin{document}

\begin{center}
{\Large\bfseries Synthesis Matrix and Organizational Outline}\\[0.35em]
{\normalsize Draft Literature Review Plan: Dr. Lowell Feedback Revision}\\[0.25em]
{\normalsize Md Kamruzzaman Kamrul}\\
{\normalsize EDCI 59100-016 \textbar{} Summer 2026}
\end{center}

\vspace{0.6em}

\section{Purpose and Submission Map}

This assignment moves the review from individual article notes to a working structure for synthesis. The review topic is the use of vision-language models and related multimodal artificial intelligence systems for construction-site safety monitoring and activity tracking.

Spatiotemporal reasoning is used here in a simple way. It means understanding what is happening, where it is happening, and how that situation changes over time. On a construction site, this may include PPE use, unsafe proximity, equipment activity, work progress, or whether a visual scene matches a safety rule.

\textbf{Review type.} The final review is planned as a scoping review using PRISMA-ScR guidance. This fits the project because the literature is still developing across construction management, computer vision, natural language processing, robotics, and safety science. The studies use different data sources, model types, and evaluation metrics. A scoping review can chart that range without forcing the papers into a narrow effect-size comparison.

\section{Source Selection}

The full working corpus includes baseline, transition, core, and review papers. For this assignment, I selected 15 key sources. The goal is not to include every paper. The goal is to choose papers that help reveal patterns, contrasts, and gaps.

The selected set includes:

\begin{itemize}
\item baseline papers that show the earlier dependence on sensors, manual reporting, and computer vision;
\item transition papers that move from detection toward semantic reasoning, captioning, and activity recognition;
\item core VLM or multimodal studies for safety monitoring and progress tracking;
\item review papers that help locate broader gaps in construction AI and computer vision.
\end{itemize}

\section{Synthesis Matrix}

The revised matrix below responds to Dr. Lowell's feedback by making the synthesis visible inside each entry. Instead of separating description, connection, and outline use, each row states how the study contributes to the field's movement from detection, to interpretation, to reasoning and interaction. The final columns identify the repeated limitation or shift across studies and translate the source into an analytical claim for the literature review.

\begin{landscape}
\footnotesize
\rowcolors{2}{softgray}{white}
\begin{longtable}{L{2.7cm} L{3.2cm} L{6.2cm} L{4.9cm} L{5.1cm}}
\caption{Revised Synthesis Matrix for Key Sources}\\
\toprule
\rowcolor{lightblue}
\textbf{Source} & \textbf{Role in field progression} & \textbf{Integrated synthesis} & \textbf{Repeating limitation or shift} & \textbf{Analytical claim for the review} \\
\midrule
\endfirsthead
\toprule
\rowcolor{lightblue}
\textbf{Source} & \textbf{Role in field progression} & \textbf{Integrated synthesis} & \textbf{Repeating limitation or shift} & \textbf{Analytical claim for the review} \\
\midrule
\endhead
\bottomrule
\endfoot

Zhong et al. (2019) & Baseline map of construction computer vision. & This review establishes that construction computer vision was already used for inspection, safety monitoring, performance analysis, and tracking before current VLM work. It therefore provides the baseline from which later studies move: early systems could detect and classify visual site information, but they rarely connected that information to construction meaning, safety reasoning, or language-based interpretation. & The field begins with task-specific visual analysis. The repeated limitation is not the absence of automation, but the limited semantic depth of that automation. & Earlier construction AI research created the visual-monitoring foundation for VLM work, but its task-specific focus left room for approaches that interpret what visual evidence means in context. \\

Sherafat et al. (2020) & Baseline for worker and equipment activity recognition. & This review shows that construction researchers had already developed strong methods for recognizing selected worker and equipment activities. It builds the temporal branch of the literature, but it also shows why later progress-reporting and VLM studies were needed: recognizing an activity is not the same as explaining its meaning, linking it to schedule progress, or producing a usable report for managers. & Earlier activity-recognition systems can perform well in defined settings, but their outputs often remain isolated labels rather than narrative or decision-ready information. & Activity recognition is an important precursor to spatiotemporal VLM research because it identifies site actions, while later studies try to turn those actions into interpretable safety and progress knowledge. \\

Nath et al. (2020) & Baseline for visible safety compliance detection. & This study demonstrates that real-time deep learning can detect workers, hardhats, and safety vests with practical safety value. It connects directly to later PPE and hazard-interpretation studies because it shows both the strength and boundary of detection: the model can identify visible compliance, but it does not determine whether a broader work situation is safe. & Across safety studies, the field repeatedly moves from visible PPE detection toward contextual hazard interpretation. & PPE detection gives the review a clear starting point for safety automation, but it also illustrates why construction safety monitoring must move beyond object labels. \\

Wu et al. (2021) & Transition from detection to semantic safety reasoning. & This paper builds on object-detection safety work by adding semantic reasoning and ontology-based representation for hazard identification. It begins to close the gap between seeing objects and interpreting safety rules, but the reasoning is still attached to the visual model rather than naturally fused within it. This makes the study a bridge between traditional CV and ontology-supported VLM systems. & The literature shifts from detecting site entities toward representing relationships among entities, rules, and hazards. & Semantic reasoning marks a middle stage in the field: construction AI starts to interpret safety meaning, but still needs stronger visual-language integration. \\

Liu et al. (2020) & Transition from visual recognition to language-based scene description. & This study moves the literature from recognizing construction activity toward describing scenes in natural language. It prepares the path for VisualSiteDiary and video-to-report systems because it shows that visual site information can become readable text. At the same time, its dependence on dataset size and image richness shows that language generation does not automatically solve the construction-domain data problem. & Captioning adds interpretation, but early captioning is mostly static and highly dependent on construction-specific training data. & Image captioning is a key bridge from detection to reporting, but it also reveals the recurring need for richer datasets and stronger temporal evidence. \\

Gil and Lee (2024) & Shift from supervised PPE detection to zero-shot visual-language safety monitoring. & This study extends PPE detection by using image captioning and visual-language matching to reduce dependence on large task-specific annotation sets. Compared with Nath et al., it shows a more flexible route to safety monitoring; compared with later rule-matching and agent systems, it remains closer to classification than full contextual reasoning. & Zero-shot flexibility reduces annotation burden, but field reliability still depends on prompt quality, image complexity, and domain vocabulary. & Zero-shot PPE monitoring reflects the field's movement from fixed supervised detection toward more adaptable visual-language safety interpretation. \\

Chen et al. (2024) & Shift from monitoring to interactive worker safety support. & This study moves safety monitoring toward interaction by combining augmented reality, deep learning, and visual safety queries. Instead of only detecting hazards for later review, the system supports real-time warning and training. This connects to assistant-style systems because it treats safety information as something workers can query and use during work. & The field shifts from passive detection to interactive feedback, but usability, latency, and field deployment remain unresolved. & Interactive safety systems show that VLM-related construction research is moving toward decision support, not only automated observation. \\

Wang et al. (2024) & Shift from visual compliance checks to visual-text safety-rule interpretation. & This study connects image content with textual safety requirements through semantic similarity. It develops the same movement signaled by Wu et al.: safety monitoring becomes stronger when visual evidence is interpreted against rules rather than treated as isolated labels. However, the value of this approach depends on whether the generated visual description and the rule language actually align. & Across safety papers, the main shift is from identifying what is visible to interpreting whether the visible condition violates a safety rule. & Visual-text rule matching strengthens the review's argument that construction safety AI is moving toward contextual hazard interpretation. \\

Chan et al. (2025) & Shift from generic VLM use to domain-guided VLM safety agents. & This study extends ontology-based safety reasoning by placing construction knowledge closer to the VLM workflow through domain-specific ontology and prompting. It responds to a repeated limitation in earlier work: generic models often lack construction context. By adding domain knowledge, the system moves from open-ended VLM interpretation toward more controlled, construction-aware reasoning. & Domain knowledge helps reduce hallucination and context gaps, but the approach still depends on prompt design, domain data, and validation. & Domain-guided VLM agents represent a later stage in the field, where reasoning is no longer only visual-language matching but construction-specific interpretation. \\

Hussain et al. (2026) & Current endpoint of assistant-supported safety inspection. & This paper frames safety inspection as an intelligent-assistant process that combines visual evidence with language-based interpretation. It brings together the interaction emphasized by Chen et al. and the domain-guided reasoning emphasized by Chan et al., showing how safety monitoring is becoming more conversational, interpretive, and decision-oriented. & Assistant systems make safety interpretation more usable, but they raise unresolved questions about cross-site reliability, accountability, and trust. & Assistant-supported inspection shows the field's movement toward interactive reasoning systems that can explain and discuss safety conditions rather than only flagging detections. \\

Bui et al. (2026) & Benchmark evidence for VLM limits in worker-state interpretation. & This study evaluates general VLMs on construction worker images and shows that even strong models struggle with close worker categories and static-image limits. It connects backward to activity-recognition research and forward to temporal VLM studies because it demonstrates that recognizing worker states from still images is not enough for dynamic construction behavior. & The limitation shifts from whether models can generate interpretations to whether those interpretations are reliable for close, dynamic, construction-specific categories. & Benchmark results show that current VLMs are promising but still weak for temporal and fine-grained worker-state reasoning in construction sites. \\

Jung et al. (2024) & Shift from scene captioning to practical daily-reporting workflows. & VisualSiteDiary extends earlier image-captioning work by turning site images into daily-reporting support and text-based retrieval. This makes the language output more practically useful than simple scene description. However, the study also keeps a recurring limitation visible: report quality depends on dataset coverage, construction vocabulary, and evaluation choices. & The progress-monitoring branch moves from describing scenes toward producing usable documentation, but it remains sensitive to data coverage and metric selection. & Vision-language reporting systems show how multimodal AI can translate site observations into management documents, while still needing stronger evaluation standards. \\

Xiao et al. (2024) & Shift from activity tracking to video-supported report generation. & This study connects earlier activity-recognition work with readable report generation by using video-based productivity analysis and language-supported reporting. It therefore shows an important progression: activity labels become part of a construction-management workflow. At the same time, the evaluation combines detection performance and human judgment, making comparison with other studies difficult. & The literature moves toward workflow integration, but evaluation metrics become more mixed as systems combine detection, generation, and user-facing reports. & Video-to-report systems show that progress monitoring is becoming more interpretive and document-oriented, but the field needs clearer benchmarks for multimodal outputs. \\

Jeoung et al. (2025) & Direct move toward spatial, task, and temporal event reasoning. & This study advances the progress-monitoring branch by separating spatial, task, and temporal event questions for equipment monitoring with MLLMs on video frames. It extends earlier activity-recognition work because the goal is no longer only to label an activity, but to answer structured questions about what is happening, where, and when. & The field begins to address temporal reasoning directly, but still lacks long-horizon and real-site temporal benchmarks. & Spatiotemporal question answering is the clearest evidence that construction VLM research is moving from recognition toward reasoning about site events over time. \\

Zhang et al. (2025) & Data-efficient recognition using pretrained vision-language methods. & This study shows how pretrained vision-language methods can reduce the need for retraining in construction object recognition. It connects to zero-shot and few-shot safety studies by addressing the same data problem from a tool/material recognition angle. Still, open-vocabulary recognition remains dependent on reliable prompts, domain terms, and construction-specific validation. & Few-shot and training-free methods reduce annotation burden, but they do not remove the need for domain-aware evaluation. & Data-efficient VLM methods help explain why the field is moving toward pretrained models, while also showing that data scarcity remains a central constraint. \\

\end{longtable}
\end{landscape}

\section{Theme Development and Categorization}

The four themes below now frame the literature as a progression rather than as separate topic bins. Each theme represents a stage in the field's development: first recognizing objects and tasks, then interpreting safety meaning, then integrating reasoning into reporting and interaction, and finally confronting the limitations that still shape the field.

\subsection{Theme 1: Stage 1 - Recognizing Objects, Tasks, and Site Activity}

\textbf{Theme definition.} The first stage of the field focuses on recognizing what is visible on a construction site: workers, equipment, PPE, tools, materials, and selected activities. This stage matters because it creates the visual-monitoring foundation for later VLM work.

\textbf{Supporting studies.} Zhong et al. (2019), Sherafat et al. (2020), Nath et al. (2020), Gil and Lee (2024), and Zhang et al. (2025).

\textbf{Synthesis.} Zhong et al. and Sherafat et al. show that construction AI had already developed substantial capacity for computer vision and activity recognition before current VLM work. Nath et al. make this foundation concrete through real-time PPE detection. Gil and Lee and Zhang et al. show that the field then begins to loosen its dependence on large task-specific datasets through zero-shot, few-shot, and pretrained visual-language methods. Across these studies, the repeated movement is from fixed detection pipelines toward more flexible recognition systems.

\textbf{Tension.} Recognition is necessary, but it is not enough. Detecting PPE, equipment, tools, or worker actions does not by itself explain whether a site condition is safe, whether progress is acceptable, or whether a visual scene matches a rule. This limitation creates the need for the second stage of interpretation.

\subsection{Theme 2: Stage 2 - Interpreting Safety Meaning and Contextual Hazards}

\textbf{Theme definition.} The second stage moves beyond visible compliance checks toward contextual safety interpretation. Here, the literature begins to connect visual evidence with safety rules, worker behavior, unsafe proximity, ontology-based reasoning, and domain-specific construction knowledge.

\textbf{Supporting studies.} Nath et al. (2020), Wu et al. (2021), Gil and Lee (2024), Wang et al. (2024), Chan et al. (2025), Hussain et al. (2026), and Bui et al. (2026).

\textbf{Synthesis.} Nath et al. and Gil and Lee establish PPE monitoring as a practical safety starting point. Wu et al. and Wang et al. show the next step: safety systems must represent relationships and compare visual evidence with safety meaning. Chan et al. adds construction ontology and prompting to make VLM interpretation more domain-aware, while Hussain et al. frames inspection as an assistant-supported process. Bui et al. adds caution by showing that even advanced VLMs can struggle with close worker-state categories in static images.

\textbf{Tension.} The same limitation repeats across this stage: safety is not only a visual property. A worker may be wearing PPE and still be unsafe because of equipment proximity, work sequence, fall risk, or site context. Later VLM-based systems try to address this issue, but they introduce new concerns about prompt reliability, hallucination, accountability, and field validation.

\subsection{Theme 3: Stage 3 - Integrating Reasoning, Reporting, and Interaction}

\textbf{Theme definition.} The third stage connects recognition and interpretation to practical construction-management outputs. The literature begins to generate captions, daily reports, visual queries, and spatiotemporal answers that can support safety inspection, progress tracking, and worker-facing interaction.

\textbf{Supporting studies.} Liu et al. (2020), Chen et al. (2024), Jung et al. (2024), Xiao et al. (2024), Jeoung et al. (2025), and Hussain et al. (2026).

\textbf{Synthesis.} Liu et al. starts the shift from recognition to language by generating construction scene captions. Jung et al. and Xiao et al. develop this into daily reporting and video-supported documentation, showing that multimodal systems can support construction-management workflows. Chen et al. and Hussain et al. extend the same movement into interaction, where users can query safety conditions rather than only receive model outputs. Jeoung et al. makes the spatiotemporal dimension most explicit by separating spatial, task, and temporal event questions for equipment monitoring.

\textbf{Tension.} As systems become more useful for reporting and interaction, their evaluation becomes harder. Detection accuracy alone is not enough for captions, reports, VQA responses, or assistant recommendations. The field therefore needs clearer ways to judge whether generated outputs are accurate, useful, timely, and trustworthy.

\subsection{Theme 4: Cross-Cutting Limits - Temporal Reasoning, Benchmarks, and Field Deployment}

\textbf{Theme definition.} The fourth theme names the limits that repeat across all stages. These include limited construction-specific datasets, inconsistent metrics, domain shift, occlusion, site noise, weak long-horizon temporal reasoning, scarce real-site validation, and lack of shared benchmarks.

\textbf{Supporting studies.} Liu et al. (2020), Zhong et al. (2019), Xiao et al. (2024), Chan et al. (2025), Bui et al. (2026), Jeoung et al. (2025), and Zhang et al. (2025).

\textbf{Synthesis.} Liu et al. shows that captioning performance depends heavily on dataset size and richness. Bui et al. shows that general VLMs struggle with fine-grained construction worker categories. Chan et al. demonstrates the need for domain knowledge to control VLM reasoning. Xiao et al. shows that evaluation becomes mixed when systems combine detection, language generation, and human judgment. Jeoung et al. shows progress toward temporal reasoning, but also makes clear that long-horizon real-site benchmarks remain limited.

\textbf{Tension.} Synthetic data, few-shot learning, zero-shot learning, and pretrained VLMs reduce annotation cost. They do not remove the need for reliable field validation. A model that works on curated images may still struggle with occlusion, poor lighting, clutter, camera angle, site-specific terminology, changing worker behavior, or long sequences of construction activity.

\section{Organizational Outline for the Literature Review}

\subsection{I. Introduction}

\begin{enumerate}
\item Open with the practical problem. Construction sites are dynamic work environments where safety and progress information changes quickly.
\item Define the review topic. The review examines VLMs and related multimodal AI systems for safety monitoring and activity tracking on active construction sites.
\item Define spatiotemporal reasoning in plain language. The review uses the term to mean understanding what is happening, where it is happening, and how it changes over time.
\item State the review purpose. The purpose is to map model architectures, data modalities, construction tasks, evaluation practices, deployment barriers, and gaps.
\item State the method. This is a scoping review using PRISMA-ScR guidance because the literature is broad, mixed, and still developing.
\end{enumerate}

\textbf{Guiding questions.}

\begin{itemize}
\item What VLM or multimodal architectures are being used for construction-site safety monitoring and activity tracking?
\item How are visual, textual, temporal, and geometric data combined to support construction safety or progress decisions?
\item What barriers appear when these models are used for real-time or field-based monitoring?
\item How are these systems evaluated, and what metrics or benchmarks are used across studies?
\end{itemize}

\subsection{II. Scoping Review Approach}

This section will explain the search and charting process. It will not read like a full systematic review methods section, but it will keep the review transparent.

\begin{enumerate}
\item Identify the databases and search process, including Web of Science, Scopus, and backward snowballing.
\item Explain inclusion and exclusion criteria. Core studies should focus on active construction sites, safety monitoring, progress tracking, activity recognition, or multimodal construction scene understanding.
\item Explain why baseline papers are included. They show the shift from manual, sensor-based, NLP-only, or CV-only systems toward multimodal reasoning.
\item Describe charting categories. These include model type, data modality, construction task, temporal component, evaluation metric, deployment context, and limitation.
\end{enumerate}

\subsection{III. Stage 1: Recognizing Objects, Tasks, and Site Activity}

This body section will make the first analytical claim: construction VLM research grows out of earlier object, PPE, equipment, and activity recognition work, but those systems mainly identify what is visible rather than explaining what the visible evidence means.

\begin{itemize}
\item Establish the baseline with Zhong et al., Sherafat et al., and Nath et al.
\item Explain how Gil and Lee and Zhang et al. shift recognition toward more flexible zero-shot, few-shot, and pretrained visual-language methods.
\item End with the synthesis claim that the field first solves parts of the recognition problem before confronting the harder problem of interpretation.
\end{itemize}

\subsection{IV. Stage 2: Interpreting Safety Meaning and Contextual Hazards}

This body section will make the second analytical claim: safety monitoring evolves from checking visible compliance to interpreting contextual hazard meaning.

\begin{itemize}
\item Begin with PPE detection as the visible-compliance baseline.
\item Compare supervised detection, zero-shot PPE monitoring, semantic reasoning, and visual-text rule matching.
\item Show how ontology-supported VLM agents and assistant systems attempt to add domain context.
\item End with the synthesis claim that practical safety judgment requires context, rules, behavior, and site conditions, not only object labels.
\end{itemize}

\subsection{V. Stage 3: Reasoning, Reporting, Interaction, and Temporal Monitoring}

This body section will make the third analytical claim: VLM-related systems become more useful when they translate visual evidence into reports, answers, and interactive support, but this shift also makes evaluation more complex.

\begin{itemize}
\item Begin with image captioning as the bridge from recognition to language.
\item Compare photolog captioning, daily reporting, and video-supported report generation.
\item Connect reporting systems with query-based and assistant-style safety systems.
\item Use Jeoung et al. to discuss spatial, task, and temporal event questions.
\item End with the synthesis claim that the field is moving toward reasoning and interaction, while long-horizon temporal monitoring remains underdeveloped.
\end{itemize}

\subsection{VI. Cross-Cutting Limits: Datasets, Metrics, and Field Deployment Barriers}

This body section will make the fourth analytical claim: the same limits repeat across stages of the literature, even as the models become more advanced.

\begin{itemize}
\item Discuss limited construction-specific datasets.
\item Discuss zero-shot, few-shot, pretrained, and ontology-supported methods as partial responses to data and context limits.
\item Compare evaluation metrics across detection, captioning, VQA, report generation, and user studies.
\item Discuss field deployment issues such as occlusion, lighting, clutter, prompt reliability, domain shift, temporal continuity, and real-time use.
\item End with the synthesis claim that the field needs shared benchmarks and stronger real-site validation before VLM systems can be trusted in routine construction monitoring.
\end{itemize}

\subsection{VII. Conclusion Placeholder}

The conclusion will summarize the main pattern. Construction AI research is moving from detecting visible site elements toward interpreting safety and progress meaning through language, context, and multimodal reasoning. The strongest future directions are construction-specific VLM benchmarks, better temporal datasets, transparent field validation, and systems that connect visual observations to safety rules and work progress in ways practitioners can trust.

\section{Submission Note}

This document contains all three required components in one formatted file: the synthesis matrix, the theme development section, and the organizational outline. If the course platform requires separate uploads, the matrix and theme sections can be exported as the working document, and the outline section can be submitted as the separate outline document.

\end{document}
